\documentclass[10pt,conference]{IEEEtran}
\usepackage[T1]{fontenc}
\usepackage{amsmath,amssymb,amsthm,graphicx,url}
\usepackage[hidelinks]{hyperref}
\newtheorem{theorem}{Theorem}
\newtheorem{lemma}{Lemma}
\newtheorem{proposition}{Proposition}
\theoremstyle{definition}
\newtheorem{definition}{Definition}

\usepackage{booktabs}

% Long tokens (tool invocations, file names) make two-column justification
% tight; a little emergency stretch is preferable to overfull lines.
\emergencystretch=2em

% Inline code. Carets, percent signs and tildes are escaped at the call site.
\newcommand{\code}[1]{\texttt{#1}}
% Guarantee names.
\newcommand{\G}[1]{G#1}

\newtheorem{remarknum}{Remark}

\begin{document}

\title{Recoverable Releases Across Package Registries:\\
       A Git-Backed Saga Protocol}

\maketitle

\begin{abstract}
Releasing interdependent packages writes irreversibly to independent
registries without a coordinator or rollback. A partial failure can therefore
leave a consumer permanently behind a dependency that was published.

This paper presents CCME, a commit-message language and release protocol for
resumable releases, and dispat, its open-source implementation for single-package
repositories, monorepos, and polyrepos joined through a control repository.
In the formal model, release state consists of git tags, the
plan is a pure function of repository history, dependency graph and
configuration, and propagated work is admitted against the consumer's pending
window. We prove eight guarantees, including no orphans, exactly-once ledger
discharge and convergence under an explicit retry invariant. Minimal counterexamples establish why delivery must be tracked per consumer,
why prerelease aggregation needs a separate window, and why retry guarantees
require stable channel and suppression predicates.

Against the official dispat~1.8.2 image, the fresh campaign passes six selected
CCME~2.0 regressions, agrees with a restricted reference oracle on all 1{,}500 core
instances, recovers 24/24 injected-failure episodes within four runs, and
completes eight scale cells. The full profile intentionally generates
regressive prereleases: 997 plans were accepted and 503 complete diagnostic
plans were rejected with E195. All 1{,}914 differing rows retain their
heuristic classification and stated coverage limits.
\end{abstract}

\begin{IEEEkeywords}
release engineering, semantic versioning, sagas, partial failure,
idempotency, deterministic replay
\end{IEEEkeywords}

\section{Introduction}
\label{sec:intro}

A release workspace can contain one package, several packages in a monorepo,
or packages drawn from separately maintained repositories (a polyrepo).
Releasing it means choosing versions for the packages that need them,
writing those versions into manifests, building artefacts,
uploading them to registries, and recording what happened. In a modern
project the registries are plural and heterogeneous: an npm package sits
next to a Go module, a Docker image, and a documentation site, all wired
into one dependency graph.

The recording is the hard part. Independent registries offer no atomic
transaction spanning the release: each publish leg may succeed or fail, and
the outcome may be lost before the corresponding git tag is written. This is
a long-running transaction over separately committed resources, the setting
for sagas~\cite{sagas}; recovery must complete or reconcile successful legs
rather than assume a cross-registry rollback.

Existing tools expose different ledgers and recovery boundaries. Lerna,
Nx~\cite{nx}, release-please, and Changesets~\cite{changesets} automate
releases, while semantic-release~\cite{semanticrelease} targets one package
per repository. Their pinned fault experiments below show that registry
queries, intent files, and tag timing answer different parts of what an
interrupted cross-registry run still owes.

This paper treats the release as the transaction it is, and makes three
contributions:
\begin{itemize}
  \item a commit-message language, CCME, in which release intent
    (blast radius on two separate, one-way-coupled axes, version bumps and
    prerelease channels) is reviewable text in the commit
    (Section~\ref{sec:ccme});
  \item a formal model of the release computation and publish protocol,
    self-contained in this paper, with eight guarantees proved over it
    (Sections~\ref{sec:model} and~\ref{sec:guarantees}); minimal counterexamples identify necessary design choices and
    the explicit invariant under which retries preserve the guarantees
    (Section~\ref{sec:eval:formal});
  \item dispat, an ecosystem-agnostic implementation in which a package is
    a directory and a stage is a shell command, evaluated by differential
    validation, selected CCME~2.0 vectors, recovery campaigns, replay, and
    pinned fault experiments (Sections~\ref{sec:impl} and~\ref{sec:eval}).
\end{itemize}

\section{Motivating Example}
\label{sec:motivation}

A workspace holds a library \code{core} and three consumers \code{cli},
\code{ui}, \code{api}, each depending on \code{core}. A commit $C$ lands
with the message \code{feat(core)\^{}:~streaming reader}: \code{core}
takes a minor bump, and the caret propagates a patch bump one edge
outward. The run publishes \code{core@1.5.0}, \code{ui@0.9.2} and
\code{api@1.2.1}, tagging each on success, and then \code{cli}'s publish
fails, for any of the reasons registry writes fail.

The natural expectation is that re-running finishes the job. Whether it
does depends on one definition. Suppose the engine keeps a commit's
release intent alive only while the commit's \emph{own} package still has
it pending, the natural reading. Then on the second run \code{core}
has released $C$, so the unit has no source packages, propagates to
nothing, and \code{cli} is never released, on this or any future run.
\code{cli} is \emph{orphaned}: permanently one patch behind a dependency
it declares, with no diagnostic, because from the engine's point of view
nothing is pending. The failure is silent, and grows more likely with
the fan-out.

The repair is a one-line change of definition: a unit propagates to a
dependent exactly while the \emph{dependent's} pending window still
contains the unit's commit (Definition~\ref{def:adm}). Catch-up then
needs no repair pass or side state when prior outcomes are confirmed by tags;
an uncertain external publish/tag gap still requires identity reconciliation.
Section~\ref{sec:guarantees} proves the guarantee this buys (\G{2}, no
orphans) and its siblings.

\section{CCME: The Language}
\label{sec:ccme}

CCME is a strict superset of Conventional Commits~1.0.0~\cite{cc} over a
stated conforming subset. One commit message is a sequence of
\emph{units} separated by a whole-line separator (\code{---}), each an
independent conventional-commit-shaped block with its own header, body
and footers, so a squash-merged pull request describes several unrelated
changes without one lie of a type for all of them:

\begin{small}
\begin{verbatim}
feat(@acme/core)^^minor%beta++*: reader
---
fix(@acme/cli): correct exit code on SIGINT
---
cancel(@acme/legacy-adapter): reset state
\end{verbatim}
\end{small}

\emph{Scopes name packages, or are derived.} The scope is a
comma-separated package set, globs with exclusions
(\code{feat(*,-docs)}), or absent, in which case the affected set is
derived from the files the commit changed, package by owning directory.
Derivation makes ordinary unscoped commits release the packages they touched.

\emph{Two axes, two depths.} A commit answers two separate questions
about its dependents: how far does the \emph{version bump} travel, and
how far does the \emph{prerelease channel} travel. Both default to depth
zero: a commit that says nothing releases its own packages and nothing
else. The bump axis is written \code{\^{}} (direct consumers),
\code{\^{}\^{}} (transitive closure) or \code{+N}; the channel axis
\code{\%} (own packages change channel), \code{\%\%} (dependents do),
\code{++N} (how far). Every sigil desugars to a footer; the header above
desugars to
\begin{small}
\begin{verbatim}
feat(@acme/core): reader

Propagate: minor
Channel: beta
Propagate-Depth: all
Propagate-Channel-Depth: all
\end{verbatim}
\end{small}
reading: \code{core} takes a minor bump and enters the beta line, every
transitive dependent takes a minor bump and follows it onto \code{beta}.
Channel values may be \emph{transitions} \code{<from>><to>}: move only
packages currently on \code{<from>}. Transitions make graduation
idempotent: \code{release(@acme/*)\%beta>stable:~ship} graduates
whatever is still on beta and can be replayed safely after a partial
failure.

\begin{table*}[t]
\caption{Header examples.}
\label{tab:headers}
\small
\begin{tabular}{@{}ll@{}}
\toprule
Header & Meaning \\
\midrule
\code{feat: add retry} & derived scope, minor bump \\
\code{fix(api): null guard} & \code{api}, patch bump \\
\code{perf(core)\^{}patch+1: x} & patch to direct consumers \\
\code{refactor(core)\^{}\^{}major!: x} & breaking; major to closure \\
\code{feat(cli)\%beta: watch mode} & \code{cli} onto the beta line \\
\code{feat(core)\^{}\%beta++1: x} & self beta; consumers follow \\
\code{release(a)\%beta>stable: ship} & graduate those still on beta \\
\code{cancel(*): reset state} & discard pending records \\
\bottomrule
\end{tabular}
\end{table*}

\emph{Cancellation and corrections.} The type \code{cancel} discards a
scope's accumulated unreleased release metadata; a unit is cancelled for
a package when its commit is an ancestor-or-self of the cancelling
commit; ancestry, unlike timestamps, is deterministic under merges and
rebases. The footers \code{Edits:~<sha>} and \code{Deletes:~<sha>}
restate or discard a named past record: a commit
\code{fix(core):~rewrite internals} carrying \code{Edits:~<sha of A>}
replaces $A$'s pending major with this unit's patch, so \code{core}
releases \code{1.4.3} rather than \code{2.0.0}. Corrections reach only
unreleased work and only strict ancestors; the newest correction wins; a
correction whose own record is discarded is void.
\code{Release-As:~none} holds a package until lifted; an exact
\code{Release-As} pins one version.

\section{The Model}
\label{sec:model}

This section defines the release computation and the publish protocol;
Section~\ref{sec:guarantees} proves the guarantees over these
definitions and nothing else.

\subsection{Workspace, State, Windows}

\begin{definition}[history]
The history is a finite DAG of commits with a distinguished
$\mathsf{HEAD} = h$; $\mathit{reach}(c)$ is the set of commits reachable
from $c$ by parent edges, including $c$; write $c' \preceq c$ for
$c' \in \mathit{reach}(c)$. $h$ is fixed throughout. This reading of a
repository is the standard one: the commit DAG with reachability is the
working model of git history in repository mining~\cite{birdgit}, and
the same structure underlies formal derivations of git's three-way
merge from least common ancestors~\cite{mrdt}. Unlike the
patch-theoretic accounts of version
control~\cite{swierstraloh,hpt,mimram}, this model tracks no file
contents at all, only history, tags and messages, which is what
makes the plan a pure function of the repository.
\end{definition}

\begin{definition}[workspace]
$\mathcal{P}$ is the finite set of packages at $h$; $E \subseteq
\mathcal{P} \times \mathcal{P}$ contains $(a,b)$ when $a$ declares a
dependency on $b$ of a configured kind; $\mathit{dep}(p) = \{d : (d,p)
\in E\}$. The graph is acyclic (a cycle is a repository-scoped error: no
publish order exists for it, so whichever member went first would
declare a range on a version that does not exist). For $S \subseteq
\mathcal{P}$, $\delta(S,d)$ is the least number of edges on a chain from
a member of $S$ to $d$ against the arrows, $\infty$ if none.
\end{definition}

\begin{definition}[bumps, versions, channels]
$\mathbb{B} = \{\mathsf{none} < \mathsf{patch} < \mathsf{minor} <
\mathsf{major}\}$ is a chain, hence a join-semilattice under
$\sqcup = \max$. Versions are SemVer~2.0.0~\cite{semver} with its
precedence; $\mathit{ch}(v)$ is $\mathsf{stable}$ for a non-prerelease
$v$ and otherwise $v$'s first prerelease identifier; $\mathit{core}(v)$
and $\mathit{cnt}(v)$ are the numeric triple and the prerelease counter;
$\mathit{apply}(v,b)$ applies a bump.
\end{definition}

\begin{definition}[state]
\label{def:state}
A state $\sigma$ is a finite set of release tags $(p, v, c)$ with
$c \preceq h$; states are ordered by inclusion. $\mathit{base}_\sigma(p)$
is the highest-precedence version among $p$'s tags,
$\mathit{sbase}_\sigma(p)$ the highest-precedence \emph{non-prerelease}
one, $\mathit{basec}_\sigma(p)$ and $\mathit{scom}_\sigma(p)$ their
commits; $\mathit{ch}(\mathit{base}_\sigma(p)) = \mathsf{stable}$ by
convention when $p$ has no tag. Tags are never moved or deleted, so
states only grow: a tag is the record that a leg of the transaction
completed, and a rewritten record is what makes a saga unrecoverable.
Tag write access is therefore release authority; a manifest's version field
is not the release ledger. Manifests, configuration, and derived scopes remain
trusted planning inputs.
\end{definition}

\begin{definition}[train and fresh-work windows]
\label{def:window}
\begin{equation}
  W_\sigma(P) \;=\; \mathit{reach}(h) \setminus
                    \mathit{reach}(\mathit{scom}_\sigma(P)),
  \label{eq:window}
\end{equation}
and
\begin{equation}
  W^\circ_\sigma(P) = W_\sigma(P) \setminus
                       \mathit{reach}(\mathit{basec}_\sigma(P)),
  \label{eq:fresh}
\end{equation}
with an empty subtracted reach when $P$ has no tag.  $W_\sigma$ is the
\emph{train window}; $W^\circ_\sigma$ is the \emph{fresh-work window}.
When $P$ has no stable baseline, $W_\sigma(P) = \mathit{reach}(h)$:
an unreleased package's first release accounts for its whole
history, which is also why adoption on an existing repository needs no
migration step. The window is measured from the last \emph{stable} tag
so that a prerelease train can recompute its target over the whole train
(Eq.~\ref{eq:train}). Direct admission, propagation admission and ledger
discharge use $W^\circ$ instead. Thus a prerelease tag discharges delivered
work without erasing the earlier records needed to recompute the train target.
\end{definition}

\subsection{Units, Suppression, Admission}

\begin{definition}[units]
$\mathcal{U}$ is the set of units carried by commits in
$\bigcup_P W_\sigma(P)$, parsed and with scopes resolved and corrections
applied. For $u \in \mathcal{U}$: $c_u$ is its commit,
$\mathit{scope}(u) \subseteq \mathcal{P}$ its resolved scope,
$\mathit{bump}(u) \in \mathbb{B}$ the bump its type and \code{!} map to;
its bump axis is $\mathit{pb}(u) \in \mathbb{B}$ (the stated value, or
$\mathit{bump}(u)$ under \code{inherit}; the default once propagation is
asked for is $\mathsf{patch}$, the honest signal for a dependent whose
own API did not change), $\mathit{pd}(u) \in \mathbb{N} \cup
\{\infty\}$ (default $0$), $\mathit{ps}(u)$ a scope filter; its channel
axis is a channel value $\mathit{pc}(u)$, a depth $\mathit{pcd}(u)$
(default $0$) and a filter $\mathit{pcs}(u)$.
\end{definition}

\begin{definition}[suppression and sources]
\label{def:src}
$\mathit{canc}(c,p)$ holds when $c$ is an ancestor-or-self of a
\code{cancel} commit whose scope contains $p$; $\mathit{held}(p)$ when
$p$'s effective \code{Release-As} is \code{none}. Both are functions of
$h$ and the configuration alone. With $\mathit{disch}_\sigma(p,c)
\equiv c \notin W^\circ_\sigma(p)$,
\begin{equation}
\begin{aligned}
  \mathit{src}_\sigma(u) = \{ p \in \mathit{scope}(u) : {}&
    \mathit{disch}_\sigma(p,c_u) \\
    &\lor\lnot(\mathit{canc}(c_u,p) \lor \mathit{held}(p)) \}.
\end{aligned}
  \label{eq:src}
\end{equation}
Suppression applies only to \emph{undischarged} work: once $p$ has
published the version carrying $u$, the artefact its consumers are owed
is public, and neither a later cancel nor a later hold may retroactively
strand them.
\end{definition}

\begin{definition}[reach]
\label{def:reach}
$R(S,\ell) = \{ d \in \mathcal{P} \setminus S : \delta(S,d) \leq \ell
\}$. Excluding all of $S$ means a unit never propagates to a package it
bumps directly; depth is shortest-path and always measured from the
originating source set, and a propagated effect never cascades its own
propagation, so the depth written in the message is the whole blast
radius in every run.
\end{definition}

\begin{definition}[channels]
\label{def:chan}
$\mathit{chan}_\sigma(P)$, resolved for \emph{every} package: the value
proposed by the newest direct \code{Channel} directive naming $P$ with
$c_u \in W^\circ_\sigma(P)$; else the value proposed for $P$ by the newest
channel-axis unit $u$ with $\mathit{pcd}(u) > 0$, $P \in
R(\mathit{src}_\sigma(u), \mathit{pcd}(u)) \cap \mathit{pcs}(u)$,
$c_u \in W^\circ_\sigma(P)$ and $\lnot\mathit{canc}(c_u,P)$; else
$\mathit{ch}(\mathit{base}_\sigma(P))$. A proposal is computed against
$\mathit{cur} = \mathit{ch}(\mathit{base}_\sigma(P))$ (the tag, never
a value computed this run) and is empty when it proposes
$\mathit{cur}$ itself; a transition proposes only where $\mathit{cur}$
matches its \code{<from>}; a \emph{propagated} non-transition
$\mathsf{stable}$ never graduates a prerelease package (graduation must
be deliberate); \code{inherit} resolves to the unit's own
\code{Channel} value if any, else to
$\mathit{ch}(\mathit{base}_\sigma(q))$ for its sources $q$, bytewise
least name deciding a disagreement.
\end{definition}

\begin{definition}[admission and effective bump]
\label{def:adm}
Let $\mathit{live}(u)$ mean $\mathit{bump}(u)\neq\mathsf{none}$,
$\mathit{pb}(u)\neq\mathsf{none}$, and $\mathit{pd}(u)>0$. With
$\mathit{srcCh}_\sigma(u)$ equal to
$\{\mathit{chan}_\sigma(p) : p \in \mathit{src}_\sigma(u)\}$ and
\begin{equation}
  \mathit{rsv}_\sigma(u,D) \equiv \mathsf{stable} \in
  \mathit{srcCh}_\sigma(u) \lor \mathit{chan}_\sigma(D) \in
  \mathit{srcCh}_\sigma(u)
  \label{eq:resolvable}
\end{equation}
(an installer selects a prerelease only on its own line, so a bump
across a channel boundary is not a real obligation),
\begin{equation}
\label{eq:admission}
\begin{aligned}
  \mathit{adm}_\sigma(u,D) \;\equiv\; & \\[-0pt]
     & D \in R(\mathit{src}_\sigma(u), \mathit{pd}(u))
     &\!\!& \text{(reached)} \\
  \wedge\;\; & D \in \mathit{ps}(u)
     &\!\!& \text{(in scope)} \\
  \wedge\;\; & c_u \in W^\circ_\sigma(D)
     &\!\!& \text{(fresh window)} \\
  \wedge\;\; & \neg\,\mathit{canc}(c_u, D)
     &\!\!& \text{(not discarded)} \\
  \wedge\;\; & \mathit{live}(u)
     &\!\!& \text{(bump-live)} \\
  \wedge\;\; & \mathit{rsv}_\sigma(u, D)
     &\!\!& \text{(resolvable),}
\end{aligned}
\end{equation}
\begin{equation}
\label{eq:effective}
  \mathit{eff}_\sigma(P) = \mathit{dir}_\sigma(P) \sqcup
  \bigsqcup \{\, \mathit{pb}(u) : \mathit{adm}_\sigma(u,P) \,\},
\end{equation}
where $\mathit{dir}_\sigma(P) = \bigsqcup\{\mathit{bump}(u) : P \in
\mathit{src}_\sigma(u) \cap \mathit{scope}(u),\ c_u \in W^\circ_\sigma(P)\}$.
Because $\sqcup$ is a join, a package with its own \code{feat} pending
is unaffected by an incoming propagated \code{patch}, propagation can
only raise a version, and gathering order cannot matter.
\end{definition}

The consumer-side staleness audit is the upward dual of
Eq.~\ref{eq:admission}, not merely reverse reachability. For candidate $D$
it scans units in $W^\circ_\sigma(D)$, excludes $D$ when it is itself a
source of that unit, and applies the same shortest depth, scope,
cancellation, $\mathit{live}$, and resolved-channel predicate $\mathit{rsv}$. These clauses
make the upward answer exactly equal to the downward planner's obligations,
including prerelease workspaces.

\subsection{Plan, Versions, Runs}

\begin{definition}[plan and versions]
\label{def:plan}
Write $\mathit{pin}_\sigma(P)=v$ when the effective surviving exact
\code{Release-As: v} unit for $P$ has $c_u\in W^\circ_\sigma(P)$; otherwise
$\mathit{pin}_\sigma(P)$ is undefined. Thus a successful tag discharges a
pin just like other fresh work, and the directive is fixed during a retry.
\begin{equation}
\begin{aligned}
  \Pi(\sigma) = \{\, P :\;
    &\lnot\mathit{held}(P) \wedge (\mathit{eff}_\sigma(P) \neq
     \mathsf{none} \\
    &\ \lor\ \mathit{chan}_\sigma(P) \neq
     \mathit{ch}(\mathit{base}_\sigma(P))
     \\
    &\ \lor\ \mathit{pin}_\sigma(P)\downarrow) \,\},
\end{aligned}
\end{equation}
each member carrying $\nu_\sigma(P)$. Let $S=\mathit{sbase}_\sigma(P)$,
$B=\mathit{base}_\sigma(P)$, and $E$ be the train aggregate over $W_\sigma(P)$.
If $\mathit{pin}_\sigma(P)\downarrow$, it supplies that exact version. A stable package that
remains stable uses $\mathit{apply}(S,\mathit{eff}_\sigma(P))$; a prerelease
package graduating to stable uses $\mathit{apply}(S,E)$ as required by
CCME~2.0 \S11.5. Otherwise, with $\mathit{ch}$ the target channel,
\begin{equation}
\label{eq:train}
\begin{aligned}
  t &= \mathit{apply}(S, E), \\
  \nu_\sigma(P) &=
  \begin{cases}
    t\text{-}\mathit{ch}.(\mathit{cnt}(B){+}1)
      & \begin{array}{l}\text{if } \mathit{core}(B) = t\text{ and}\\
        \mathit{ch}(B) = \mathit{ch},\end{array}\\
    t\text{-}\mathit{ch}.0 & \text{otherwise,}
  \end{cases}
\end{aligned}
\end{equation}
with one guard: every computed or pinned version must exceed
$\mathit{base}_\sigma(P)$ in precedence, or the run refuses. Recomputing
$t$ from the \emph{stable} baseline every run is what moves a whole
train to \code{2.0.0-beta.0} when a breaking change lands mid-train
(Table~\ref{tab:train}).

\end{definition}

\begin{table}[t]
\caption{Eq.~\ref{eq:train} tracing one train.}
\label{tab:train}
\small
\begin{tabular}{@{}lll@{}}
\toprule
Baseline & Pending bump, channel & Next \\
\midrule
\code{1.2.3}        & minor, \code{beta}          & \code{1.3.0-beta.0} \\
\code{1.3.0-beta.0} & no fresh commit             & empty plan \\
\code{1.3.0-beta.0} & a fresh minor commit        & \code{1.3.0-beta.1} \\
\code{1.3.0-beta.1} & minor $+$ a new fix         & \code{1.3.0-beta.2} \\
\code{1.3.0-beta.2} & a breaking change lands     & \code{2.0.0-beta.0} \\
\code{1.3.0-beta.2} & none, switch to \code{rc}   & \code{1.3.0-rc.0} \\
\code{1.3.0-rc.1}   & none, \code{stable}         & \code{1.3.0} \\
\bottomrule
\end{tabular}
\end{table}

\begin{definition}[run]
\label{def:run}
A run from $\sigma$ is parameterised by an adversarially chosen set
$F \subseteq \Pi(\sigma)$ of failing publishes, which induces the
blocked set
$B = \{p \in \Pi(\sigma) \setminus F : p$ transitively depends on some
$q \in F\}$, computed over the \emph{full} workspace graph (a package
with no bump this run is still a path from a dependent to a failed
dependency). Publication proceeds in a topological order of the full
graph, ties broken bytewise by name, filtered to the plan \emph{last}:
inducing the graph on the plan first deletes ordering paths that run
through unplanned packages. Each package is tagged immediately after its
own publish succeeds, never before (a tag for an unpublished version
makes it unrecoverable, since Eq.~\ref{eq:window} treats it as
released), never batched (a mid-run failure could repeat an external write,
which a target may reject or may perform twice), and \emph{whatever its publish target},
including targets that upload nothing
(Proposition~\ref{prop:tagall}). The run yields
$\mathit{run}(\sigma,F) = \sigma \cup \{(p, \nu_\sigma(p), h) : p \in
\Pi(\sigma) \setminus (F \cup B)\}$. A \emph{run sequence} is
$\sigma_1 \leq \sigma_2 \leq \cdots$ with $\sigma_{k+1} =
\mathit{run}(\sigma_k, F_k)$.
\end{definition}

For scale, let $H,A$ count reachable commits and parent edges, $C$ commits in
the union window, $U$ units, $I$ resolved unit--package incidences, and $Z$
retained or emitted contribution/provenance incidences. Window classes can
share immutable reachability sets, but constructing them costs in $H+A$, not
only $C$; stable and fresh windows require different boundary keys.
Literal propagation costs $O(U(|\mathcal P|+|E|)+Z)$. A graph walk may be
cached safely when source set, depth and edge kinds match; merging by window and
directive spelling is unsafe because admission, cancellation, scope,
self-source exclusion, inherited channel and provenance remain per unit.
Lexical parsing may be cached by repository, immutable record, parser/schema
and parsing-configuration identities. Scope and changed-path resolution has a
distinct workspace/path/graph key and must be recomputed or cached separately.
Channel incidence inversion costs $O(U+I+|\mathcal P|)$. These are informative
implementation costs, not required end-to-end bounds; semantic equivalence and
deterministic output, including provenance, remain normative. 

Order matters for a reason tags cannot express: a dependent published
before its dependency carries a range reconciled against the old
version, both land at the same commit, and no later run can detect the
inversion from ancestry, so it is prevented, not diagnosed. When a
publish is rejected because the version already exists (a run died
between publish and tag), CCME~2.0 requires content-identity verification
before adoption and otherwise aborts: tagging a foreign artefact would enrol
it into the release lineage. This is a protocol obligation, not a claim that
dispat can inspect arbitrary user-defined publish scripts; such integrations
must supply an idempotency or identity check.

\begin{theorem}[determinism]
\label{thm:det}
$\Pi(\sigma)$, $\nu_\sigma$, the publish order and the diagnostics are
functions of $(\sigma, h, \mathit{config})$ alone.
\end{theorem}
\begin{proof}
Every definition above is a closed expression in those inputs; the three
choice points (competing directives, competing propagated channels,
publish-order ties) are each resolved by a stated total order: newest
commit then last unit within it, and bytewise least name.
\end{proof}

\subsection{Instantiation: What the Theorems Say About a Repository}
\label{sec:model:inst}

The definitions above are parametric, and a theorem about the model
is worth nothing to an operator unless it is a theorem about
\emph{their} repository. A concrete repository at a commit determines
one \emph{instance}
\[
  I \;=\; (\text{commits and parent edges},\ h,\
           \mathcal{P}, E,\ \mathcal{U},\ \sigma,\ \mathit{config}):
\]
the commit DAG and \code{HEAD} are read from the repository, the
packages and edges from its manifests and configuration, the units
from its commit messages, and $\sigma$ from its tags. Every quantity
the model defines is a computable function of $I$: the sets are
finite, ancestry is decidable, and no definition consults anything
outside $I$. The model is a fragment of repository behaviour: it abstracts
artefact contents, hook semantics, and configuration features not represented
in $I$. The theorems of Section~\ref{sec:guarantees} are universally
quantified over instances of that fragment, subject to acyclicity, the
version guard of Definition~\ref{def:plan}, and the retry invariant stated
below. The latter is a property of the run sequence, checked again before
each retry, rather than a static property of one instance. What remains between the theorems and a running binary is a
\emph{refinement} claim, that the implementation computes this
model's $\Pi(\sigma)$ and $\nu_\sigma$ and executes
Definition~\ref{def:run}'s protocol. That claim is not proved here;
it is tested, by Section~\ref{sec:eval:fresh}'s vectors (each
fixes a small instance and the plan the model assigns it) and by
Section~\ref{sec:eval:diff} on randomised instances.

\section{Guarantees, with Proofs}
\label{sec:guarantees}

Cross-run claims use the \emph{retry invariant}: the corrected tuple
stream, source sets, graph, scope/depth predicates, target baselines, and
resolved source/target channel admission of every outstanding unit remain
unchanged, except for baseline tags written for successful packages.
Two recurring components deserve names. \textbf{(H1)} says every baseline
channel read by \code{Propagate-Channel: inherit} remains unchanged;
\textbf{(H2)} says no contribution initially suppressed by cancellation,
hold, or correction is later released at the same $h$. H1 and H2 are
necessary in their respective cases but are not, by themselves, the full
retry invariant. Neither is needed for one invocation's determinism.

\begin{lemma}[window antitonicity]
\label{lem:w}
For $\sigma' = \mathit{run}(\sigma,F)$, $W^\circ_{\sigma'}(p)=\emptyset$
if $p$ published and $W^\circ_{\sigma'}(p)=W^\circ_\sigma(p)$ otherwise.
The train window $W$ also empties after a stable release, but may remain
unchanged after a prerelease release.
\end{lemma}
\begin{proof}
The run adds only tags at $h$. If $p$ gained none, both baseline commits
are unchanged. If it published, $\mathit{basec}_{\sigma'}(p)=h$, so
Eq.~\ref{eq:fresh} subtracts $\mathit{reach}(h)$ and is empty. Only a
stable release also makes $\mathit{scom}_{\sigma'}(p)=h$.
\end{proof}

\begin{lemma}[source stability]
\label{lem:s}
Under the retry invariant, $\mathit{src}_{\sigma_j}(u)$ is constant for
every outstanding unit throughout a valid run sequence.
\end{lemma}
\begin{proof}
The corrected tuple stream and its cancellation, hold, correction, and
source-set predicates are fixed components of the invariant. H2 is the
component excluding later release of an initially suppressed contribution.
Arbitrary tag-set inclusion outside a valid run sequence implies no such
monotonicity, because a higher-version tag may point to an earlier commit.
\end{proof}

\begin{lemma}[retry stability]
\label{lem:c}
Under the retry invariant, an unreleased package's resolved channel,
each outstanding unit's source-channel set, and
$\mathit{rsv}(u,d)$ remain constant.
\end{lemma}
\begin{proof}
These are precisely the channel and source-set components fixed by the
retry invariant. Definition~\ref{def:chan} then selects the same direct
or propagated proposal, and Eq.~\ref{eq:resolvable} reads the same sets.
\end{proof}

\begin{remarknum}
\label{rem:drift}
(H1) cannot be replaced by requiring stable completed releases: \code{inherit} reads
$\mathit{ch}(\mathit{base}_\sigma(q))$, which the origin's own
graduation moves: with \code{core} at baseline \code{1.5.0-beta.0} and
commits \code{feat(core)\^{}++1:~x} and
\code{release(core)\%stable:~graduate}, run 1 plans the consumer at
\code{2.0.1-beta.0} (inherited \code{beta}) while \code{core}
graduates; if the consumer fails, run 2 plans it at \code{2.0.1}:
same package, both plans, different versions, every completed release
stable. Thus unchanged inherited-channel state must be an explicit component of the
retry invariant; successful stable releases alone do not imply version stability.
\end{remarknum}

\begin{theorem}[the eight guarantees]
\label{thm:g}
Write $\mathit{adm}_j$ for $\mathit{adm}_{\sigma_j}$ and $n =
|\Pi(\sigma_1)|$. Under the retry invariant, for every unit $u$, consumer
$D$ and run $j \geq 1$:
\begin{description}
  \item[\G{1}] \textbf{termination.} propagation halts after $O(|\mathcal{U}|
    \cdot (|\mathcal{P}|{+}|E|))$ work, whatever the declared depths.
  \item[\G{2}] \textbf{no orphans.} $\mathit{adm}_1(u,D) \Rightarrow
    \mathit{adm}_j(u,D)$ until $D$ releases at a commit containing
    $c_u$.
  \item[\G{3}] \textbf{version stability.} if $D \in \Pi(\sigma_1)$ and $D$ is not
    released before run $k$, then $D \in \Pi(\sigma_k)$ and $\nu_k(D) =
    \nu_1(D)$.
  \item[\G{4}] \textbf{no double ledger delivery.} once $D$ releases past $c_u$,
    $\mathit{adm}_j(u,D)$ never holds again; with \G{2}, ledger discharge is
    exactly-once for every contribution whose target eventually succeeds.
    This does not make external side effects exactly-once.
  \item[\G{5}] \textbf{no widening.} $\{D : \mathit{adm}_j(u,D)\} \subseteq
    \{D : \mathit{adm}_1(u,D)\}$.
  \item[\G{6}] \textbf{convergence.} $\Pi(\sigma_{n+1}) = \emptyset$, provided
    each nonempty retry admits and publishes at least one initially
    publishable, non-held package.
  \item[\G{7}] \textbf{channel discharge.} a propagated channel contribution is
    admitted only while its commit remains in $W^\circ(D)$; baseline
    equality and transition matching independently remove no-op proposals.
  \item[\G{8}] \textbf{one-way axis independence.} suppressing a bump never
    suppresses a channel change. The converse is false: resolved channels
    are an input to bump admission in Eq.~\ref{eq:resolvable}.
\end{description}
\end{theorem}

\begin{proof}
\G{1}: the traversal computing $R$ marks each package seen at most
once, so it does $O(|\mathcal{P}|{+}|E|)$ work per unit regardless of
depth or cycles; acyclicity is needed for the publish order, not here.
\G{2}: compare Eq.~\ref{eq:admission} clause by clause. Source set,
depth, graph, scope, cancellation and channel admission are fixed by the
retry invariant. The consumer's fresh window is unchanged until $D$
releases (Lemma~\ref{lem:w}); that release is the only event that can
make its admission clause false.
\G{3}: $D$ gained no tag, so its baselines, windows and
$\mathit{dir}$ are unchanged; the propagated join is unchanged by \G{2}
($\geq$) and \G{5} ($\leq$), whose proofs below use only the lemmas; so
$\mathit{eff}$ is unchanged and the channel is unchanged
(Lemma~\ref{lem:c}). Thus bump-only and channel-only membership remain;
an exact \code{Release-As} remains the same because its unit is fixed by the
invariant. In every case $\nu$ is the same function of unchanged inputs.
\G{4}: by Lemma~\ref{lem:w}, $W^\circ(D)$ empties after either a stable
or prerelease tag, killing fresh direct and propagated admission. This
also makes $\mathit{pin}_\sigma(D)$ undefined. This is discharge in the tag
ledger; an interruption after an external publish
but before its tag is handled by content-identity reconciliation, not by
an exactly-once claim about the external operation.
\G{5}: every clause of Eq.~\ref{eq:admission} is constant except the
window, which is antitone (Lemma~\ref{lem:w}); a conjunction of
constant and antitone predicates is antitone.
\G{6}: $\Pi(\sigma_{j+1}) \subseteq \Pi(\sigma_j)$, since a member of
the later plan was not released (else \G{4}) and each disjunct of
$\Pi$ is then unchanged or smaller; by the premise each nonempty retry
removes at least one initially publishable package (\G{4}), so $|\Pi|$
strictly decreases to $0$ within
$n$ runs. A run that publishes nothing has not failed \emph{partially}
and is surfaced as a hard error, not retried.
\G{7}: the channel-axis admission test uses $c_u\in W^\circ(D)$, so
Lemma~\ref{lem:w} discharges it. While fresh, equality with the baseline
and a transition's \code{<from>} check separately reject no-ops.
\G{8}: the channel pass runs first and reads no bump-axis output; the bump
pass then reads its resolved channels through Eq.~\ref{eq:resolvable}.
\end{proof}

\begin{proposition}[universal tagging is necessary]
\label{prop:tagall}
If a released package with no artefact to upload and a non-empty effective
bump were versioned but not tagged, its window would never advance
(Lemma~\ref{lem:w}), it would re-enter every plan forever, and \G{6} would
fail. Hence
Definition~\ref{def:run} tags whatever is released, and the
convergence check (re-plan after a fully successful run; a non-empty
result is an error) admits exactly one exclusion, held packages.
\end{proposition}

Three corollaries carry the operational weight. \emph{Resumption}: after
a partial failure, $\Pi(\sigma_2) = F \cup B$ at the versions of run 1
(\G{2}--\G{5}), so the next run does the rest of the same job with no
flags and no cleanup. \emph{Idempotency}: running twice with no intervening
commits, tag changes, or configuration changes gives the same plan
(Theorem~\ref{thm:det});
running after full success gives the empty plan (\G{4}).
\emph{Review soundness}: while the retry invariant holds, by \G{5} the plan a reviewer approves before
the first attempt is an upper bound on everything every retry can
publish, so an attacker who can induce a publish failure cannot enlarge
what a commit releases, and no resume needs re-approval.

\begin{table*}[t]
\caption{Section~\ref{sec:motivation} under a schedule that loses
\code{cli} twice. \G{2}: \code{cli} is in every plan until it
publishes. \G{3}: at \code{2.0.1} in all of them. \G{4}: the three that
published leave. \G{5}: nothing new enters. \G{6}: the plan reaches
empty. Run 2 publishes nothing, which \G{6} treats as the hard-error
case rather than the partial-failure one.}
\label{tab:runs}
\small
\begin{tabular}{@{}llll@{}}
\toprule
Run & Plan & Published & Failed \\
\midrule
1 & \code{core}, \code{cli}, \code{ui}, \code{api}
  & \code{core} 1.5.0, \code{ui} 0.9.2, \code{api} 1.2.1 & \code{cli} \\
2 & \code{cli} 2.0.1 & --- & \code{cli} \\
3 & \code{cli} 2.0.1 & \code{cli} 2.0.1 & --- \\
4 & $\emptyset$ & --- & --- \\
\bottomrule
\end{tabular}
\end{table*}

\section{Release-Tool Comparison}
\label{sec:compare}

In the pinned configurations we executed, Lerna and nx derive bumps from
commits and tag before publish, so a failed upload is absent from their tag
ledger; release-please has the same ordering in its documented workflow.
Changesets tags after publish, consumes intent files, and queries npm during
recovery.
CCME tags each package after its publish leg and therefore retains debt in git,
subject to the external publish/tag reconciliation limit of
Definition~\ref{def:run}. The observed propagation policies were unconditional
for Lerna/nx, range-gated for Changesets, and per-commit for CCME.
Table~\ref{tab:compare}
summarises the resulting operational distinctions; the artifact contains the
full command transcripts.
The refusal and push fixtures pin Lerna~10.0.1, Nx~23.1.2, and
Changesets~3.0.1. Lerna/nx tagged the denied version before its failed upload;
Changesets left it untagged and recovered by querying npm. In the concurrent
push fixture, compared runs left tags outside branch ancestry, while dispat
merged before its final push and planned the colleague's commit next.

\begin{table}[t]
\caption{Pinned release-tool configurations.}
\label{tab:compare}
\scriptsize
\resizebox{\columnwidth}{!}{%
\begin{tabular}{@{}llll@{}}
\toprule
 & lerna / nx & changesets & CCME / dispat \\
\midrule
bump source & commits & intent files & commits \\
tag time & before publish & after publish & after each leg \\
propagation & unconditional & range-gated & per commit \\
recovery ledger & tags / CI & npm & git tags \\
external gap & manual/query & npm query & reconcile identity \\
\bottomrule
\end{tabular}}
\end{table}

\section{Implementation}
\label{sec:impl}

dispat is a statically linked Go binary whose release-planning state is only
configuration, history, and tags. Packages are folders and lifecycle stages
are shell commands; scanners and byte-preserving writers cover 35 manifest
formats in twenty ecosystems. The planner implements the definitions above;
topological execution continues independent subtrees after a failure and
skips consumers of failed packages. A remote reserved tag serialises release
runs. Successful external effects are not rolled back: later tag, record, or
push failure leaves reconciliation work under the publish-then-record protocol. Invoked tools,
including Terraform, retain responsibility for their operational state and
identity checks. The implementation is tested through vectors, black-box CLI
runs, fuzzing, real manifests, and the campaigns below.


\paragraph{Repository layouts and saga granularity.}
For a single-package repository, $\mathcal{P}=\{p\}$ and $E=\varnothing$:
propagation is vacuous, but version planning, publish-then-tag recording and
recovery still apply. This is the one-leg case of the saga protocol. Dispat
requires the package in a subfolder; the configuration and release ledger
remain at the root. Multiple destinations inside one publish script still
share one completion tag, so partial external success requires reconciliation;
separate package legs provide separate completion records. A single source
project can expose its CLI, container image and documentation as such legs.
For a polyrepo, a control repository pins linked repositories as git submodules
inside package wrapper folders. Its pointer-update commits express release
intent; its history and tags provide the model's single $h$ and $\sigma$.
The linked repositories retain their own histories. This supplies one release
graph without merging the source repositories, rather than a distributed
transaction over independent git histories. This construction establishes
applicability; the evaluation does not measure polyrepo operational scale.

\paragraph{Reusable automation.}
A Windows developer and a Linux CI runner can share \code{global.yaml}
(graph, scripts and lifecycle). Each platform file references it with
\code{\$ref: ./global.yaml}; \code{linux.yaml} sets
\code{shell: ["/bin/sh", "-c"]}, while \code{windows.yaml} sets
\code{["powershell.exe", "-NoProfile", "-Command"]} as its shell.
Only platform overrides are repeated; references are file-relative and
replace whole keys, not nested fields. Preview with
\code{dispat --config windows.yaml status} locally or
\code{dispat --config linux.yaml status} on the runner; CI runs the latter
configuration without \code{status} after its release gates.
\code{dispat if OS=Windows\_NT} can select these commands through
\code{--then} and \code{--else}, but 1.8.2 executes those branches through
\code{/bin/sh} before the child reads its configuration. Native Windows can
select its file directly; the conditional wrapper still needs a POSIX shell.
Named scripts can themselves use \code{\$ref};
\code{dispat for Debug Release --do "dispat exec build-one"}
passes each value through \code{DISPAT\_ITEM}.

\paragraph{Agent release gates.}
A lock-free \code{status --require-release --strict} check can precede
an all-package workspace suite with verified temporary-link cleanup.
Package checks then validate reconciled versions and lockfiles, followed by
smoke tests of the exact artifacts before publication. Wrappers preserve
failures, and release authorization remains separate.
\section{Evaluation}
\label{sec:eval}

The primary campaign uses the official dispat~1.8.2 image and implemented
CCME~2.0 semantics~\cite{ccme2}. Release-tool comparisons retain their pinned
1.7.x measurements and are not represented as new 1.8.2 experiments. The
artifact records images, seeds, raw streams, checksums, and harness versions.

\subsection{Differential Validation of the Refinement Claim}
\label{sec:eval:diff}

Section~\ref{sec:model:inst} left one obligation open: that the binary
computes this paper's model. We test it with selected specification vectors and a restricted reference
oracle in approximately 300 lines of Python. The oracle predates the complete
model presented here: agreement within its core domain is evidence for that
fragment, not a proof of refinement of the full model. A seeded generator
produces instances with 5--20 packages, layered DAGs, 4--12 units, and mixed
tag states. Each instance becomes a real Git repository and is planned by
the oracle and the released binary. Golden cases gate the comparison.

The core profile isolates the oracle's shared semantic domain. The full
profile deliberately exercises prerelease baselines, channel-only entries,
and bump-less directive types outside that domain. Differences in the latter
are classified diagnostically; they are not counted as passing conformance
checks. The current campaign below reports both profiles, including rejected
plans and the limits of its heuristic classification.

\subsection{Released-Binary Evaluation}
\label{sec:eval:fresh}

We reran released CCME~2.0.0 against the official dispat~1.8.2 ARM64 image.
The campaign identity binds source tag commit \code{a7011cbcb9e8}, image digest
\code{sha256:4c0ad4b8...e7bfd}, the verified release binary, and 753 artifacts.

The revised runner validates Git fixtures and CLI exits, uses tmpfs for E1,
and distinguishes semantic refusals from execution failures. These timings
are not compared with runs using an earlier harness. All 1{,}500 core
instances agreed, with no invalid sample. Of 1{,}500 full-profile histories, which include
regressive prereleases, 997 produced accepted plans (546 clean comparisons)
and 503 produced complete E195 diagnostic refusals; four refused dossiers had
no row difference, giving 550 diagnostic-clean instances overall. Every
proposed regression set in the 503 dossiers was independently checked. The
remaining 1{,}914 package-row differences retained PG1 862, PG2 375, PG2x 14,
and PG3? 663, with UNKNOWN 0. Rejected diagnostic plans are neither executable
release plans nor conformance results. The CLI also
passed selected 2.0 vectors 133--138 (fresh discharge, suppression, channel
eligibility/order, inherited-channel drift, and source exclusion). Vector 135
observes downward W208, not the upward API; six cases do not prove refinement.

Two orphan repetitions recovered the same version to an empty third plan;
two campaigns converged in 24/24 episodes within four runs with zero structural
G3/G5 errors. Five-release replay reproduced 36/42 google-cloud-node and 7/64
conventional-changelog versions, without implying those releases used CCME.
Three fresh dispat fault cells passed 8 orphan, 12 clean-push, and 14
conflicting-push checks; competitor cells were not rerun.

All eight timing cells are valid contended-host observations under a two-CPU,
4-GiB cap, not cross-tool comparisons. Medians span 0.017\,s for 100 layered
packages to 19.354\,s for a 10{,}000-package chain; observed maximum RSS spans
15--118\,MiB. U1000 took 0.290\,s/26\,MiB and 5{,}000 padding commits took
0.664\,s/30\,MiB. The recorded fixtures passed \code{git fsck}, and the
\code{commit-tree} generator matched the ordinary small fixture's commit IDs,
refs, messages, and log output.

\subsection{Defect-Finding Power of the Formalisation}
\label{sec:eval:formal}

Minimal counterexamples identify why several apparently simpler designs fail.
They serve as regression obligations and explain the model's assumptions.

\begin{enumerate}
  \item \emph{Aggregation is not discharge.} Reusing a stable-baseline
    train window for fresh admission can produce \code{beta.1},
    \code{beta.2}, $\ldots$ without new commits. Separate $W$ and
    $W^\circ$ retain train intent while discharging delivered work.
  \item \emph{Suppression must remain stable across retries.} If a
    cancelled or corrected source is later released, its discharged record
    can change the source set and reach new consumers. H2 excludes this
    widening within a reviewed retry sequence; changing it requires a new plan.
  \item \emph{Inherited channels are mutable inputs.} The graduation
    example in Remark~\ref{rem:drift} changes a failed consumer's next version
    unless H1 preserves its inherited-channel inputs.
  \item \emph{Reverse audits need the full admission predicate.} A
    consumer-side staleness audit that omits channel eligibility or target
    self-exclusion can disagree with forward propagation.
  \item \emph{Axis independence is one-way.} Resolving channels can make a
    dependency ineligible and suppress a bump; suppressing a bump does not
    suppress a channel change. Resolving channels first preserves this order.
  \item \emph{Shared graph traversal does not merge unit identity.}
    Subtracting a union of source sets excludes legitimate targets of other
    units; admitting every traversed node can include a unit's own source.
    Traversal may be shared, but admission and provenance remain per unit.
\end{enumerate}

These counterexamples motivated the released rules used in the model. Their
role here is to establish design necessities, not to claim defects in the
current specification. The selected released-binary vectors cover six cases;
they do not establish full implementation refinement.

\section{Threats to Validity}
\label{sec:threats}

\emph{Internal.} The vectors, the tests and the formal model were
produced by the same author as the system, so they share assumptions
and their agreement cannot catch an error common to all three.
The guarantees are proved of the model, not of the Go code; the link
is the vector suite, the tests and the differential run of
Section~\ref{sec:eval:diff}, which covers the planner but not the
run's protocol. The concurrent-push experiment is the strongest
evidence for that protocol, and it is one deterministic interleaving
per tool, not a sweep over every point a push could land.
\emph{External.} The self-hosting study is one repository of moderate
size administered by the tool's author; it demonstrates feasibility
and recovery on real history, not usability at organisational scale.
Ecosystem generality is by construction plus tested writers and manifest
fixtures, chosen by us rather than a field study. \emph{Construct.}
Conformance is measured against a
contract the implementation co-evolved with;
Section~\ref{sec:eval:formal} shows this cuts both ways. No user study
supports reviewability claims; we claim only that the blast radius
is written in the message, not that authors write it correctly. The
characterisation of release-please is read from its documentation;
Lerna, nx and Changesets are executed on two faults in the pinned
six-package fixture summarised in Section~\ref{sec:compare}.
The check counts recorded for those tools describe what each leaves
behind against the questions dispat's protocol answers, not a score
on a shared scale. The experiments run in containers on one machine
and one CI provider; the replay configurations for the subject
repositories were reverse-engineered from their layouts and may
misread conventions their tools resolve differently.

\section{Related Work}
\label{sec:related}

Adams and McIntosh survey release engineering as a first-class
subject~\cite{adams}; Potvin and Levenberg document Google's monorepo
practice~\cite{potvin}, whose build layer solves ordering and caching
but not versioned publication to external registries.
Among release tools, semantic-release~\cite{semanticrelease}
established commit-driven versioning; Lerna~\cite{lerna},
changesets~\cite{changesets}, Nx~\cite{nx} and
release-please~\cite{releaseplease} extend automation to monorepos
within, predominantly, one ecosystem each; Section~\ref{sec:compare}
states their computations as instances of this paper's model and
locates the structural difference in the tag discipline. Empirical studies
show declared version semantics and actual compatibility diverge
~\cite{raemaekers,bogart}; Decan et al. study dependency-network evolution
across ecosystems~\cite{decan}. Together they motivate explicit release
intent over a dependency graph. Formal treatments of
version control (semantics of VCS operations~\cite{swierstraloh},
patch theories as higher inductive types~\cite{hpt}, merges as
pushouts~\cite{mimram}) model contents and merging; this paper's
model keeps only the DAG, tags and messages, the objects repository
mining treats as git's model~\cite{birdgit} and mergeable-type
derivations build on~\cite{mrdt}. Packages and
dependency graphs as formal objects go back to the EDOS work on
distribution consistency~\cite{edos}. Within the literature surveyed here,
we found no earlier composition of these objects into this release protocol.
The protocol adapts sagas~\cite{sagas} to legs for which no acceptable
compensating action is available; we therefore choose forward recovery in
light of transaction limits described by Gray~\cite{gray}. Helland's arguments for
independently committed, idempotently reprocessable
steps~\cite{helland07,helland12} describe the discipline
\G{2}--\G{6} impose. Kahn's algorithm~\cite{kahn} underlies the
publish order; Mokhov et al.'s build-system taxonomy~\cite{mokhov}
clarifies what dispat is not: it has no incremental rebuild layer,
because unchanged packages are absent from the plan rather than
cached. Fuzzing is reported under the cautions of Klees et
al.~\cite{klees}. The nearest formal-guarantees work in operations is
Anvil~\cite{anvil}, which verifies that Kubernetes controllers
reconcile to a desired state; convergence here instead comes from a ledger,
replay, and explicit reconciliation of uncertain external outcomes.

\section{Conclusion}
\label{sec:conclusion}

Across single-package repositories, monorepos and control-repository polyrepos,
CCME treats release publication as a saga. It records intent in commit
messages, completed publication legs in tags, and admits propagated work
against each consumer's own fresh-work window. Under the retry invariant, its eight
guarantees make deterministic replay a correct recovery procedure.
Minimal counterexamples show why separate admission windows, stable retry
predicates, and per-unit provenance are necessary. The pinned release-tool fault experiments support the operational distinction in the pinned configurations:
dispat resumes a refused upload and preserves tag reachability across a
concurrent push; the compared runs required external recovery.

Remaining work is to mechanise the model and proofs in Lean or
TLA$^{+}$~\cite{lamport} and extend differential validation from plans to
complete runs, including randomised interleavings of tagging and merging.

\section*{Acknowledgments: AI Assistance}
OpenAI Codex agents assisted with drafting and revising the exposition and
proof arguments throughout Sections~\ref{sec:intro}--\ref{sec:conclusion},
writing experiment harness code, running experiments, and summarising the
resulting records. This assistance extended beyond language editing.
Reported measurements were obtained from program executions rather than
language-model predictions. The human author remains accountable for the
claims, proofs, references, and experimental interpretation.

\section*{Data Availability}

Released source and specifications are available at
an author-identifying public repository, omitted for anonymous review.
A local supplementary bundle retains proof details, vectors, fixtures,
harnesses, checksums, and raw records. A curated review bundle contains experiment scripts, selected raw and summary
records, and hashes of the original files. It is not linked in this initial
submission: the complete Git provenance and execution records expose author
and environment information, and anonymized distribution is still being
arranged. The review bundle can be provided through the program chairs.
Binaries, container images and complete Git histories remain outside that
non-executable draft. The complete archival replication package is planned
for publication with the accepted paper.

% acmart balances the final bibliography with a 1.166pt box-rounding excess.
% The rendered content fits; tolerate only this sub-2pt vertical warning.
\vfuzz=2pt
\begin{thebibliography}{21}

\bibitem{sagas}
H.~Garcia-Molina and K.~Salem.
Sagas.
In \emph{Proceedings of the ACM SIGMOD International Conference on
Management of Data}, pages 249--259, 1987.

\bibitem{gray}
J.~Gray.
The transaction concept: Virtues and limitations.
In \emph{Proceedings of the 7th International Conference on Very Large
Data Bases (VLDB)}, pages 144--154, 1981.

\bibitem{helland07}
P.~Helland.
Life beyond distributed transactions: an apostate's opinion.
In \emph{Proceedings of the 3rd Biennial Conference on Innovative Data
Systems Research (CIDR)}, 2007.

\bibitem{helland12}
P.~Helland.
Idempotence is not a medical condition.
\emph{Communications of the ACM}, 55(5):56--65, 2012.

\bibitem{kahn}
A.~B. Kahn.
Topological sorting of large networks.
\emph{Communications of the ACM}, 5(11):558--562, 1962.

\bibitem{semver}
T.~Preston-Werner.
Semantic Versioning 2.0.0.
\url{https://semver.org/spec/v2.0.0.html}, accessed 2026.

\bibitem{cc}
Conventional Commits 1.0.0.
\url{https://www.conventionalcommits.org/en/v1.0.0/}, accessed 2026.

\bibitem{ccme2}
CCME Specification 2.0.0.
Public source URL omitted for anonymous review,
accessed 2026.

\bibitem{birdgit}
C.~Bird, P.~C. Rigby, E.~T. Barr, D.~J. Hamilton, D.~M. Germ\'an, and
P.~Devanbu.
The promises and perils of mining Git.
In \emph{Proceedings of the 6th IEEE International Working Conference
on Mining Software Repositories (MSR)}, 2009.

\bibitem{mrdt}
G.~Kaki, S.~Priya, K.~C. Sivaramakrishnan, and S.~Jagannathan.
Mergeable replicated data types.
\emph{Proceedings of the ACM on Programming Languages},
3(OOPSLA), Article 154, 2019.

\bibitem{swierstraloh}
W.~Swierstra and A.~L\"oh.
The semantics of version control.
In \emph{Proceedings of the ACM International Symposium on New Ideas,
New Paradigms, and Reflections on Programming \& Software (Onward!)},
pages 43--54, 2014.

\bibitem{hpt}
C.~Angiuli, E.~Morehouse, D.~R. Licata, and R.~Harper.
Homotopical patch theory.
In \emph{Proceedings of the 19th ACM SIGPLAN International Conference
on Functional Programming (ICFP)}, 2014.

\bibitem{mimram}
S.~Mimram and C.~Di Giusto.
A categorical theory of patches.
\emph{Electronic Notes in Theoretical Computer Science}, 298, 2013.

\bibitem{edos}
F.~Mancinelli, J.~Boender, R.~Di Cosmo, J.~Vouillon, B.~Durak,
X.~Leroy, and R.~Treinen.
Managing the complexity of large free and open source package-based
software distributions.
In \emph{Proceedings of the 21st IEEE/ACM International Conference on
Automated Software Engineering (ASE)}, 2006.

\bibitem{lerna}
Lerna: a fast, modern build system for managing and publishing multiple
JavaScript packages from the same repository.
\url{https://lerna.js.org/}, accessed 2026.

\bibitem{changesets}
Changesets: a way to manage your versioning and changelogs.
\url{https://github.com/changesets/changesets}, accessed 2026.

\bibitem{semanticrelease}
semantic-release: fully automated version management and package
publishing.
\url{https://github.com/semantic-release/semantic-release},
accessed 2026.

\bibitem{nx}
Nx: smart monorepos, fast CI.
\url{https://nx.dev/}, accessed 2026.

\bibitem{releaseplease}
release-please: automated releases based on conventional commits.
\url{https://github.com/googleapis/release-please}, accessed 2026.

\bibitem{adams}
B.~Adams and S.~McIntosh.
Modern release engineering in a nutshell: Why researchers should care.
In \emph{Proceedings of the 23rd IEEE International Conference on
Software Analysis, Evolution, and Reengineering (SANER)}, 2016.

\bibitem{potvin}
R.~Potvin and J.~Levenberg.
Why Google stores billions of lines of code in a single repository.
\emph{Communications of the ACM}, 59(7):78--87, 2016.

\bibitem{raemaekers}
S.~Raemaekers, A.~van Deursen, and J.~Visser.
Semantic versioning versus breaking changes: A study of the Maven
repository.
In \emph{Proceedings of the 14th IEEE International Working Conference
on Source Code Analysis and Manipulation (SCAM)}, 2014.

\bibitem{bogart}
C.~Bogart, C.~K\"astner, J.~Herbsleb, and F.~Thung.
How to break an API: Cost negotiation and community values in three
software ecosystems.
In \emph{Proceedings of the 24th ACM SIGSOFT International Symposium on
Foundations of Software Engineering (FSE)}, 2016.

\bibitem{decan}
A.~Decan, T.~Mens, and P.~Grosjean.
An empirical comparison of dependency network evolution in seven
software packaging ecosystems.
\emph{Empirical Software Engineering}, 24(1):381--416, 2019.

\bibitem{mokhov}
A.~Mokhov, N.~Mitchell, and S.~Peyton Jones.
Build systems \`a la carte.
\emph{Proceedings of the ACM on Programming Languages},
2(ICFP):79:1--79:29, 2018.

\bibitem{klees}
G.~Klees, A.~Ruef, B.~Cooper, S.~Wei, and M.~Hicks.
Evaluating fuzz testing.
In \emph{Proceedings of the ACM SIGSAC Conference on Computer and
Communications Security (CCS)}, pages 2123--2138, 2018.

\bibitem{anvil}
X.~Sun, W.~Ma, J.~T. Gu, Z.~Ma, T.~Chajed, J.~Howell, A.~Lattuada,
O.~Padon, L.~Suresh, A.~Szekeres, and T.~Xu.
Anvil: Verifying liveness of cluster management controllers.
In \emph{Proceedings of the 18th USENIX Symposium on Operating Systems
Design and Implementation (OSDI)}, 2024.

\bibitem{lamport}
L.~Lamport.
\emph{Specifying Systems: The TLA\textsuperscript{+} Language and
Tools for Hardware and Software Engineers}.
Addison-Wesley, 2002.


\end{thebibliography}

\end{document}
